\section{Principle III: Order}

\begin{quotex}
Although he described himself as a pagan, the influence of Thomism on \textbf{Charles Maurras} is quite clear. Being and Order, for him, are coterminous. I don't know the extent of Donoso Cortes' influence on Maurras, but Maurras also rejects endless discussion, claiming that to debate order, which has a mathematical-like certainty, is a waste of time. To choose disorder (or chaos) is to choose death and annihilation. Man, such as he is, needs to hold firm to his fixed reference points (or road signs?); modernity, on the other hand, wants to knock over every one of them. Who is left to still remember them and celebrate them on the calendar? 

\end{quotex}

As no face could exist without the features that surround it and the line which contains it, as soon as the Being begins to move away from its opposite, as soon as the Being is, it has its form, it has its order, and it is that same thing from which it is limited that forms it. What existence is without essence? What is Being without law? In all the degrees of the scale, the Being weakens when order diminishes; it dissolves as soon as order no longer holds.

Order is only a means. It is a point of departure. To reestablish order restores an atmosphere favorable to the action of the spirit as well as that of the body. That order makes work possible or better. It guarantees it duration, supplies it with assistants or protectors.

Humane rule does not consist in killing, destroying, nor annihilating the subject whom it must, on the contrary, develop while maintaining him on his path.

The necessity to \emph{subordinate} in order to coordinate and order, there is no rhetorical nonsense that can go against that mathematics!

To be conformed to order shortens and facilitates work. To contradict or debate order is to waste time.

Order, they say, is a higher justice.

For historical and political order, \emph{to have} it is nothing, \emph{to hold} it is nothing, if one is not also in a position \emph{to keep} it.

In war as in peace, order is precious among all the goods. With its fake hardness, with its apparent strictness, it saves lives, as it measures and makes use of efforts.

The soldier who complains about the order is the enemy of himself. The blind kindness that joins in with this soldier is an enemy of the soldier. An unconscious and involuntary enemy; what does the intention matter if it sends him to his death?

Precisely because he is ungrateful and weak, because forgetfulness and fickleness are normal for him, man notices early on that it is necessary for him to look again, in Time that is constantly changing, for some points of fixed reference, for invariable points of support, whenever he wants to realize a plan of any importance, or he wants to be faithful to his goal and his love.

I deliberately write this last word, which only expresses the sentiment of persons; for if strong passions have their anniversary rites, if a return of certain dates leads to a natural return of thought to the joyful or sorrowful mysteries of the life of the heart, or even more so, when it is no longer a question of a sole being, but of a society, a religion, a cause, then will it be necessary to immortalize happy or gloomy memories on their dates.

\textbf{Do not forget}: this is the starting point of all order and all law.



\flrightit{Posted on 2012-10-14 by Cologero }

\begin{center}* * *\end{center}

\begin{footnotesize}\begin{sffamily}



\texttt{Jason-Adam on 2012-10-16 at 00:10 said: }

Indeed, order is the basis of all life. Without a basis there is no society, as we see in our days as all forms of collective identity are in decline.


\hfill

\texttt{Cologero on 2012-10-16 at 17:08 said: }

It is important to not that, for Maurras, order not only determines what a being is, but also what it is not. This is important for his subsequent analysis.


\hfill

\texttt{Lucius Cornelius Sulla on 2017-05-22 at 18:47 said: }

How is Charles Maurras pagan? He was pagan after he lost faith, but he was then involeved in politics it is only after Dreyfus affair, that he become engaged in political affairs, but I never i read anywere that he considers himself pagan when he was defending church and become monarchist, or can some point me to were I can find that sentence?


\hfill

\texttt{Cologero on 2017-05-22 at 19:09 said: }

Obviously, dude, you read it here. Unless you consider here to be Nowhere\footnote{\url{https://www.youtube.com/watch?v=pesIGuV9DDk}}.


\hfill

\texttt{Lucius Cornelius Sulla on 2017-05-25 at 19:17 said: }

you misinterpreted my question: how charles maurras is pagan?.It is last thing I could call him, he clear would describe himself as positivist and classicist.


\hfill

\texttt{Cologero on 2017-05-25 at 21:12 said: }

OK, Sulla, from the OED definition of pagan\footnote{\url{https://en.oxforddictionaries.com/definition/pagan}}, the most reliable source:

\emph{A person holding religious beliefs other than those of the main world religions.}

So that is the sense I used the word. You are correct; he would have described himself as a positivist, at least after his youth. His first book, Anthinéa, about his visit to Athens, was originally titled ``Pagan promenades". Eventually, however, he embraced Positivism and achieved a rapprochement with Catholicism.

At that time, he became coy about the designation ``pagan" or even ``atheist", but that need not concern us. Yet if he was a classicist, then he was a pagan in the same sense that Plato was a pagan.

Since the book under discussion was written much later in his life, I will change to text to read: ``he was often described as a pagan", or something like that.

But I hope that excessive pedantry does not obscure the main point. To wit, the recognition of a cosmic order, i.e., ``natural law", is possible for pagans, so it is not solely a Christian dogma.


\end{sffamily}\end{footnotesize}
